\documentclass[11pt]{amsart}
\usepackage{amssymb,amsmath,amsthm,mathtools}
\usepackage[margin=1.15in]{geometry}
\usepackage[colorlinks,citecolor=blue,linkcolor=blue]{hyperref}

\newtheorem{theorem}{Theorem}[section]
\newtheorem{lemma}[theorem]{Lemma}
\newtheorem{proposition}[theorem]{Proposition}
\newtheorem{corollary}[theorem]{Corollary}
\theoremstyle{definition}
\newtheorem{definition}[theorem]{Definition}
\newtheorem{example}[theorem]{Example}
\theoremstyle{remark}
\newtheorem{remark}[theorem]{Remark}

\newcommand{\Z}{\mathbb{Z}}
\newcommand{\Q}{\mathbb{Q}}
\newcommand{\C}{\mathbb{C}}
\newcommand{\HH}{\mathfrak{H}}
\newcommand{\Zbar}{\overline{\Z}}
\newcommand{\Qbar}{\overline{\Q}}
\newcommand{\SL}{\mathrm{SL}}
\newcommand{\GL}{\mathrm{GL}}
\newcommand{\im}{\operatorname{Im}}
\newcommand{\sz}{\sigma_0}
\newcommand{\Xs}{X_0(N)^*}
\newcommand{\Gs}{\Gamma_0(N)^*}
\newcommand{\smat}[4]{\left(\begin{smallmatrix}#1&#2\\#3&#4\end{smallmatrix}\right)}
\newcommand{\etanorm}[1]{\lVert\eta\rVert(#1)}

\title[Integrality of CM values on $X_0(N)^*$]{Integrality of normalized CM values of cusp forms on $X_0(N)^*$, and rigorous identification of CM points on canonical models}
\date{November 2026. Draft for internal checking.}

\begin{document}
\begin{abstract}
Let $N$ be squarefree and let $F$ be a weight-$2$ cusp form on $\Gamma_0(N)$ with
integer Fourier coefficients that is invariant under all Atkin--Lehner
involutions.  We show that
$F(\tau)^{6\sz(N)}/\prod_{d\mid N}\Delta(d\tau)$ is an algebraic integer at every
CM point $\tau$, lying in the field of definition of the image of $\tau$ on
$\Xs$, and we prove the analogous statement at elliptic points.  Two consequences
make numerical evaluation at CM points rigorous: an explicit bound for the height
of a rational CM point on the canonical model of $\Xs$, and a test that proves a
form vanishes at a CM point from a crude numerical upper bound.  The argument is
the classical one for class invariants; we make no claim of novelty, and record
the statements in the form needed for computation.
\end{abstract}
\maketitle

\section{Introduction}

Throughout, $N\ge1$ is a squarefree integer, $\sz(N)$ is its number of divisors,
and $\Gs$ is the group generated by $\Gamma_0(N)$ and its Atkin--Lehner
involutions, with quotient curve $\Xs=X_0(N)/W(N)$, a curve over $\Q$ with a single
cusp.  Let $S=S_2(\Gamma_0(N))^{W(N)}$ be the space of weight-$2$ cusp forms
invariant under all Atkin--Lehner involutions; it is identified with the
holomorphic differentials on $\Xs$, and a basis $f_1,\dots,f_g$ of $S$ with
integer Fourier coefficients defines the canonical map
$\varphi\colon\Xs\to\mathbb P^{g-1}$ over $\Q$ (for $g\ge2$).

A non-cuspidal point $x\in\Xs$ that is the image of a CM point $\tau\in\HH$ has
coordinates $\varphi(x)=(f_1(\tau):\dots:f_g(\tau))$ away from elliptic points,
and these can be computed numerically from $q$-expansions.  The problem is to
turn such numerics into proofs: to prove that a given rational point $P$ of the
canonical model \emph{is} $\varphi(x)$, or that a given hyperplane contains
$\varphi(x)$.  A numerical value alone cannot prove that a rational number
equals a given one without a bound on its height.  The results below supply
that bound.

Put
\[
D_N(\tau)=\prod_{d\mid N}\Delta(d\tau),\qquad
\etanorm{z}=\lvert\eta(z)\rvert\,\im(z)^{1/4},\qquad
T(\tau)=\prod_{d\mid N}\etanorm{d\tau}^{4/\sz(N)} ,
\]
where $\Delta=\eta^{24}$.  The function $\etanorm{\cdot}$ is
$\SL_2(\Z)$-invariant.

\begin{theorem}\label{thm:main-intro}
Let $F\in S$ have integer Fourier coefficients, let $\tau\in\HH$ be a CM point
that is not an elliptic point of $\Gs$, and let $x\in\Xs$ be its image.  Then
\[
n_F(\tau)=\frac{F(\tau)^{6\sz(N)}}{D_N(\tau)}
\]
is an algebraic integer in the residue field $\Q(x)$, and
$\lvert n_F(\tau)\rvert=\nu_F(\tau)^{6\sz(N)}$ where
\[
\nu_F(\tau)=\frac{\sqrt N\,\im(\tau)\,\lvert F(\tau)\rvert}{T(\tau)} .
\]
In particular, if $x\in\Xs(\Q)$ then $n_F(\tau)\in\Z$, and
$\nu_F(\tau)<1$ implies $F(\tau)=0$.
\end{theorem}

\begin{corollary}\label{cor:height-intro}
In the situation of Theorem~\ref{thm:main-intro}, suppose $x\in\Xs(\Q)$ and write
$\varphi(x)=(a_1:\dots:a_g)$ with coprime integers $a_i$.  Then
$\lvert a_i\rvert\le\nu_{f_i}(\tau)$ for all $i$.
\end{corollary}

Section~\ref{sec:int} proves a more general statement
(Theorem~\ref{thm:integral}); Section~\ref{sec:cons} gives the versions at
elliptic points and for points of higher degree; Section~\ref{sec:tail} gives an
explicit bound for Fourier coefficients, which is what a rigorous evaluation of
$F(\tau)$ needs; Section~\ref{sec:ex} works three examples at level $67$ by
hand; Section~\ref{sec:num} records numerical checks.

The proof is the standard argument that a modular function which is holomorphic
on $\HH$ and has integral expansions at all cusps is integral over $\Z[j]$,
combined with the integrality of singular moduli.  The first statement is
\cite[Ch.~2, \S2, Lemma~2.1]{KubertLang}; the second is
\cite[Ch.~5, \S2, Thm.~4]{Lang}, \cite[Thm.~11.1]{Cox}.  For squarefree level
and functions whose images under all Atkin--Lehner involutions have integral
expansions at $\infty$, the resulting integrality of CM values is
\cite[Lemmas~4.3 and~4.4]{BruinierOno}, which gives the integrality in
Theorem~\ref{thm:main-intro} once the function $F^{6\sz(N)}/D_N$ is written
down.  What is added here is the choice of $D_N$, the formula for the absolute
value and the resulting height bound and vanishing test, and the version at
elliptic points, where the relevant form is not Atkin--Lehner invariant.  For
functions whose divisor
is supported on CM points and cusps far more is known, namely the prime
factorization of the values \cite{GrossZagier}; the forms $F$ here have zeros at
arbitrary points, so only integrality is available, and it is all that is needed.

\section{Integrality}\label{sec:int}

For $\gamma=\smat abcd\in\GL_2^+(\Q)$ and a function $f$ on $\HH$ let
$(f|_k\gamma)(\tau)=\det(\gamma)^{k/2}(c\tau+d)^{-k}f(\gamma\tau)$; scalar
matrices act trivially.  For $Q\mid N$ an \emph{Atkin--Lehner matrix} is an
integer matrix $w_Q=\smat{Qx}{y}{Nz}{Qt}$ of determinant $Q$.  These form a single
coset $\Gamma_0(N)w_Q$ and normalize $\Gamma_0(N)$ \cite[\S2]{AL}; they
generate $\Gs$ together with $\Gamma_0(N)$.  Write $q^{1/Q}=e^{2\pi i\tau/Q}$,
$\zeta_Q=e^{2\pi i/Q}$, and let $\Zbar$ be the ring of algebraic integers.

\begin{lemma}\label{lem:coset}
Let $\gamma=\smat abcd\in\SL_2(\Z)$ and $Q=N/\gcd(c,N)$.  There are $j\in\Z$ and
an Atkin--Lehner matrix $w_Q$ with
$\gamma=w_Q\cdot\smat 1j0Q$ up to a scalar.
\end{lemma}
\begin{proof}
Since $N$ is squarefree, $\gcd(c,Q)=1$; choose $j$ with $cj\equiv d\pmod Q$.
Then $\gamma\smat Q{-j}01=\smat{Qa}{b-aj}{Qc}{d-cj}$ has determinant $Q$, its
lower left entry is divisible by $N$ and its diagonal entries by $Q$, so it is
an Atkin--Lehner matrix $w_Q$; and $\smat Q{-j}01^{-1}=Q^{-1}\smat 1j0Q$.
\end{proof}

\begin{definition}
A form $H\in M_k(\Gamma_0(N))$ is \emph{cusp-integral} if for every
$\gamma\in\SL_2(\Z)$, with $Q=N/\gcd(c,N)$,
\[
Q^{k/2}\cdot(H|_k\gamma)\in\Zbar[[q^{1/Q}]] .
\]
\end{definition}

\begin{lemma}\label{lem:AL-ci}
Let $H\in M_k(\Gamma_0(N))$ have Fourier coefficients in $\Z$ and satisfy
$H|_kw_Q=\pm H$ for all $Q\mid N$.  Then $H$ is cusp-integral.
\end{lemma}
\begin{proof}
By Lemma~\ref{lem:coset},
$H|_k\gamma=\pm H|_k\smat 1j0Q=\pm Q^{-k/2}H\bigl((\tau+j)/Q\bigr)$, and
$H((\tau+j)/Q)=\sum_n a_n\zeta_Q^{nj}q^{n/Q}$.
\end{proof}

\begin{lemma}\label{lem:delta}
Let $\Delta_d(\tau)=\Delta(d\tau)$ for $d\mid N$.
\begin{enumerate}
\item For $\gamma=\smat ab{c}{\delta}\in\SL_2(\Z)$ with $Q=N/\gcd(c,N)$ there are integers
$A,B,D'$ with $AD'=d$, $D'=\gcd(d,Q)$, such that
$(\Delta_d|_{12}\gamma)(\tau)=D'^{-12}\,\Delta\bigl((A\tau+B)/D'\bigr)$.
\item $D_N|_{12\sz(N)}\gamma=Q^{-6\sz(N)}\cdot u\cdot q^{r}\,(1+\cdots)$ with $u$
a root of unity, $r\in\frac1Q\Z_{>0}$ and $1+\cdots\in\Z[\zeta_Q][[q^{1/Q}]]$.
\item $D_N|_{12\sz(N)}w_Q=D_N$ for every Atkin--Lehner matrix $w_Q$.
\end{enumerate}
\end{lemma}
\begin{proof}
(1) Write $\smat d001\gamma=\gamma'\smat AB0{D'}$ with $\gamma'\in\SL_2(\Z)$ and
$A,D'>0$ (Hermite normal form); then $A$ is the gcd of the first column
$(da,c)$, so $A=\gcd(d,c)=\gcd(d,N/Q)$ and $D'=d/A=\gcd(d,Q)$, using that $d\mid N$
and $N$ is squarefree.  With $z=(A\tau+B)/D'$ we have $d\gamma\tau=\gamma'z$ and
$c\tau+\delta=D'\cdot j(\gamma',z)$, where $j(\gamma',z)$ is the automorphy factor,
because the bottom row of $\smat d001$ is $(0,1)$.  Hence
$(c\tau+\delta)^{-12}\Delta(d\gamma\tau)=D'^{-12}j(\gamma',z)^{-12}\Delta(\gamma'z)=D'^{-12}\Delta(z)$.

(2) The number of $d\mid N$ with $\gcd(d,Q)=e$ is $\sz(N/Q)$ for each $e\mid Q$,
and $\prod_{e\mid Q}e=Q^{\sz(Q)/2}$, so
$\prod_{d\mid N}\gcd(d,Q)^{-12}=Q^{-6\sz(Q)\sz(N/Q)}=Q^{-6\sz(N)}$.  Each factor
$\Delta((A\tau+B)/D')=\zeta q^{A/D'}\prod_{n\ge1}(1-\zeta^nq^{An/D'})^{24}$ with
$\zeta=e^{2\pi iB/D'}$.

(3) Let $w_Q=\smat{Qx}{y}{Nz}{Qt}$ and $d=d_1d_2$ with $d_1\mid Q$,
$d_2\mid N/Q$; put $d^*=(Q/d_1)d_2$.  A direct computation gives
\[
\smat d001w_Q\smat{d^*}001^{-1}
=d_1\begin{pmatrix}d_1x&d_2y\\ \frac{N}{Qd_2}z&\frac{Q}{d_1}t\end{pmatrix},
\]
and the matrix on the right is integral of determinant $1$.  Since
$\Delta_d=d^{-6}\Delta|_{12}\smat d001$, this gives
$\Delta_d|_{12}w_Q=(d^*/d)^{6}\Delta_{d^*}$.  As $d\mapsto d^*$ is a bijection of
the divisors of $N$, the product over $d$ is $D_N$.
\end{proof}

\begin{theorem}\label{thm:integral}
Let $H\in M_k(\Gamma_0(N))$ be cusp-integral and let $a,b\ge1$ be integers with
$ak=12\sz(N)b$.  Then
\[
G=H^a/D_N^{\,b}
\]
is a modular function for $\Gamma_0(N)$, holomorphic on $\HH$, and integral over
$\Zbar[j]$.  Consequently $G(\tau)$ is an algebraic integer for every CM point
$\tau\in\HH$.  Moreover, for every $\tau\in\HH$,
\begin{equation}\label{eq:nu}
\lvert G(\tau)\rvert=\nu_H(\tau)^a,\qquad
\nu_H(\tau)=\Bigl(\frac{\sqrt N\,\im\tau}{T(\tau)}\Bigr)^{k/2}\lvert H(\tau)\rvert .
\end{equation}
\end{theorem}
\begin{proof}
$G$ has weight $0$ for $\Gamma_0(N)$ and is holomorphic on $\HH$ because $\Delta$
does not vanish there.  Let $\gamma\in\SL_2(\Z)$ and $Q=N/\gcd(c,N)$.  By
cusp-integrality $H^a|\gamma\in Q^{-ak/2}\Zbar[[q^{1/Q}]]$, and by
Lemma~\ref{lem:delta}(2)
$D_N^{\,b}|\gamma=Q^{-6\sz(N)b}u^bq^{rb}(1+\cdots)$.  Since
$ak/2=6\sz(N)b$ the powers of $Q$ cancel, and $G\circ\gamma\in\Zbar((q^{1/Q}))$.

Let $\gamma_1,\dots,\gamma_\mu$ represent $\Gamma_0(N)\backslash\SL_2(\Z)$ and
$P(X)=\prod_i(X-G\circ\gamma_i)$.  Its coefficients are holomorphic on $\HH$,
invariant under $\SL_2(\Z)$, and have expansions in $\Zbar((q))$ (invariance under
$\tau\mapsto\tau+1$ removes fractional powers).  Such a function is a polynomial
in $j$ with coefficients in $\Zbar$: if its expansion begins $cq^{-m}$ with
$m>0$ then $c\in\Zbar$ and subtracting $cj^m$ lowers $m$, because
$j=q^{-1}+744+\cdots\in q^{-1}\Z[[q]]$; and a holomorphic $\SL_2(\Z)$-invariant
function bounded at $\infty$ is constant.  So $P\in\Zbar[j][X]$ is monic with
$P(G)=0$.  For CM $\tau$ the value $j(\tau)$ is an algebraic integer
\cite[Thm.~11.1]{Cox}, \cite[II, Thm.~6.1]{SilvermanAT}, hence so is $G(\tau)$.

For \eqref{eq:nu}, note that
$\lvert\Delta(d\tau)\rvert=\etanorm{d\tau}^{24}(d\,\im\tau)^{-6}$ and
$\prod_{d\mid N}d=N^{\sz(N)/2}$, so
\[
\lvert D_N(\tau)\rvert=\Bigl(\frac{T(\tau)}{\sqrt N\,\im\tau}\Bigr)^{6\sz(N)},
\]
and $6\sz(N)b=ak/2$.
\end{proof}

\begin{proposition}[Field of definition]\label{prop:field}
Let $H$, $a$, $b$, $G$ be as in Theorem~\ref{thm:integral}, let $\tau\in\HH$ have
image $x\in\Xs$, and suppose $H$ has Fourier coefficients in $\Q$.  If
$H^a|_{ak}w_Q=H^a$ for all $Q\mid N$, then $G\in\Q(\Xs)$, so $G(\tau)\in\Q(x)$ and
$\sigma(G(x))=G(\sigma x)$ for $\sigma\in\mathrm{Gal}(\Qbar/\Q)$.
\end{proposition}
\begin{proof}
By Lemma~\ref{lem:delta}(3), $G$ is invariant under $\Gs$, and its $q$-expansion
at $\infty$ is rational.  The function field of $X_0(N)$ over $\Q$ consists of
the modular functions for $\Gamma_0(N)$ with rational $q$-expansion at $\infty$
\cite[Ch.~7]{DS}, \cite[Ch.~6]{Shimura}, and the Atkin--Lehner involutions
are defined over $\Q$, so $G\in\Q(\Xs)$.  It is regular at $x$.
\end{proof}

\begin{proof}[Proof of Theorem~\textup{\ref{thm:main-intro}} and Corollary~\textup{\ref{cor:height-intro}}]
\leavevmode\\
Take $H=F$, $k=2$, $a=6\sz(N)$, $b=1$; $H$ is cusp-integral by
Lemma~\ref{lem:AL-ci} and $F^a$ is Atkin--Lehner invariant, so
Theorem~\ref{thm:integral} and Proposition~\ref{prop:field} apply.  A rational
algebraic integer is an integer, and a nonzero integer has absolute value at
least $1$.

For the corollary, $\tau$ is not elliptic, so $\tau$ is a local coordinate at $x$
and $f_i(\tau)=\lambda a_i$ for some $\lambda\in\C^\times$ (the canonical system
has no base points).  Then $n_{f_i}(\tau)=c\,a_i^{6\sz(N)}$ with
$c=\lambda^{6\sz(N)}/D_N(\tau)$, and $c\in\Q$ since some $a_i\ne0$.  For every
prime $p$ some $a_i$ is a $p$-adic unit, so $c\in\Z$, and $c\ne0$.  Hence
$\lvert a_i\rvert^{6\sz(N)}\le\lvert n_{f_i}(\tau)\rvert=\nu_{f_i}(\tau)^{6\sz(N)}$.
\end{proof}

\section{Elliptic points and points of higher degree}\label{sec:cons}

Let $\theta=q\,\frac{d}{dq}=\frac1{2\pi i}\frac{d}{d\tau}$ and let
$\vartheta_kf=\theta f-\frac k{12}E_2f$ be the Serre derivative, which maps
$M_k(\Gamma)$ to $M_{k+2}(\Gamma)$ for every finite index subgroup
$\Gamma\subset\SL_2(\Z)$ and satisfies
$\vartheta_k(f|_k\gamma)=(\vartheta_kf)|_{k+2}\gamma$ for $\gamma\in\SL_2(\Z)$.
(In \cite[\S5.1]{Zagier} this is stated for $\SL_2(\Z)$; the equivariance, and
with it the statement for subgroups, follows from the same computation, since
$\theta(f|_k\gamma)-(\theta f)|_{k+2}\gamma=\frac{k}{12}\bigl(E_2-E_2|_2\gamma\bigr)f|_k\gamma$
by the transformation law of $E_2$.)

\begin{lemma}\label{lem:serre}
If $H\in M_k(\Gamma_0(N))$ is cusp-integral then so is $12\,\vartheta_kH$.
\end{lemma}
\begin{proof}
$12\,\vartheta_k(H|\gamma)=12\,\theta(H|\gamma)-kE_2\cdot(H|\gamma)$.  On
$Q^{-k/2}\Zbar[[q^{1/Q}]]$ the operator $\theta$ multiplies the coefficient of
$q^{n/Q}$ by $n/Q$, and $E_2\in\Z[[q]]$, so the result lies in
$Q^{-(k+2)/2}\Zbar[[q^{1/Q}]]$.
\end{proof}

Let $\tau_0\in\HH$ and let $e$ be the order of its stabilizer in $\Gs$ modulo
scalars.  In the coordinate $u=(\tau-\tau_0)/(\tau-\bar\tau_0)$ the stabilizer
acts by $u\mapsto\zeta u$ with $\zeta$ a primitive $e$-th root of unity, and
$t=u^e$ is a local coordinate at the image $x$.  For $F\in S$ the differential
$F\,d\tau=g(u)\,du$ is invariant, so $g(u)=\sum_mc_mu^m$ with
$m\equiv-1\pmod e$; thus $F$ vanishes at $\tau_0$ to order at least $e-1$, and
$F\,d\tau=(c_{e-1}/e+O(t))\,dt$.  Consequently
\begin{equation}\label{eq:phi-ell}
\varphi(x)=\bigl(\theta^{e-1}f_1(\tau_0):\dots:\theta^{e-1}f_g(\tau_0)\bigr).
\end{equation}

\begin{theorem}\label{thm:elliptic}
Let $F\in S$ have integer Fourier coefficients, let $\tau_0$ be a CM point with
stabilizer of order $e$ as above, and let $x$ be its image in $\Xs$.  Put
$F^{[e]}(\tau_0)=12^{\,e-1}\,\theta^{e-1}F(\tau_0)$.  Then
\[
n_F(\tau_0)=\frac{F^{[e]}(\tau_0)^{6\sz(N)}}{D_N(\tau_0)^{e}}
\]
is an algebraic integer in $\Q(x)$, of absolute value $\nu_F(\tau_0)^{6\sz(N)}$
where
\[
\nu_F(\tau_0)=\Bigl(\frac{\sqrt N\,\im\tau_0}{T(\tau_0)}\Bigr)^{e}\,\lvert F^{[e]}(\tau_0)\rvert .
\]
For $e=1$ this is Theorem~\ref{thm:main-intro}.
\end{theorem}
\begin{proof}
Let $H=12^{e-1}\vartheta_{2e-2}\circ\cdots\circ\vartheta_2F\in M_{2e}(\Gamma_0(N))$.
It is cusp-integral by Lemmas \ref{lem:AL-ci} and \ref{lem:serre}.  By induction,
$\vartheta_{2m}\circ\cdots\circ\vartheta_2F=\theta^mF+\sum_{i<m}c_i\,\theta^iF$
with $c_i$ holomorphic on $\HH$; since $\theta^iF(\tau_0)=0$ for $i<e-1$ we get
$H(\tau_0)=F^{[e]}(\tau_0)$.  Theorem~\ref{thm:integral} with $k=2e$,
$a=6\sz(N)$, $b=e$ shows that $n_F(\tau_0)=G(\tau_0)$ is an algebraic integer of
the stated absolute value, where $G=H^a/D_N^{\,e}\in\Q(X_0(N))$.

It remains to show $G(\tau_0)\in\Q(x)$.  Let $w=\smat{*}{*}{c}{d}$ be an
Atkin--Lehner matrix of determinant $Q$, so that
$F(w\tau)=\rho(\tau)F(\tau)$ with $\rho(\tau)=(c\tau+d)^2/Q$ and
$w'(\tau)=\rho(\tau)^{-1}$.  The stabilizer of $w\tau_0$ is conjugate to that of
$\tau_0$, so $F$ vanishes to order at least $e-1$ at both points.
Differentiating $e-1$ times at $\tau_0$ and using this vanishing gives
$F^{(e-1)}(w\tau_0)=\rho(\tau_0)^{e}F^{(e-1)}(\tau_0)$.  By
Lemma~\ref{lem:delta}(3), $D_N(w\tau_0)=\rho(\tau_0)^{6\sz(N)}D_N(\tau_0)$.
The identity $H=12^{e-1}\theta^{e-1}F$ holds at $w\tau_0$ as well, and therefore $G(w\tau_0)=G(\tau_0)$.
Now let $y\in X_0(N)$ be the image of $\tau_0$.  An element of
$\mathrm{Gal}(\Qbar/\Q(x))$ maps $y$ to a point of the fibre over $x$, that is, to
the image of some $w\tau_0$, so it fixes $G(y)$.
\end{proof}

\begin{corollary}[Height bound]\label{cor:height}
Let $x\in\Xs(\Q)$ be a non-cuspidal CM point, $\tau_0$ a point above it with
stabilizer of order $e$, and $\varphi(x)=(a_1:\dots:a_g)$ with coprime integers.
Then $\lvert a_i\rvert\le\nu_{f_i}(\tau_0)$, with $\nu$ as in
Theorem~\ref{thm:elliptic}.
\end{corollary}
\begin{proof}
By \eqref{eq:phi-ell}, $f_i^{[e]}(\tau_0)=\lambda a_i$, and the proof of
Corollary~\ref{cor:height-intro} applies with
$c=\lambda^{6\sz(N)}/D_N(\tau_0)^e$.
\end{proof}

\begin{corollary}[Vanishing test]\label{cor:vanish}
Let $F=\sum_im_if_i$ with $m_i\in\Z$, and let $x\in\Xs$ be a non-cuspidal CM
point of degree $h=[\Q(x):\Q]$ with conjugates $x=x_1,\dots,x_h$, represented by
$\tau_1,\dots,\tau_h\in\HH$.  If
\[
\prod_{j=1}^h\nu_F(\tau_j)<1
\]
then $\sum_im_iX_i$ vanishes at $\varphi(x_j)$ for every $j$.  In particular this
holds if $\nu_F(\tau_1)\cdot M^{h-1}<1$, where $M\ge1$ bounds $\nu_F$ on a finite
Galois-stable set of CM points containing $x$.
\end{corollary}
\begin{proof}
The order $e$ is the same for all $x_j$.  Indeed, if $y_j\in X_0(N)$ lies above
$x_j$, then $e$ is the ramification index of $\HH\to X(1)$ at $\tau_j$, divided
by that of $X_0(N)\to X(1)$ at $y_j$, times that of $X_0(N)\to\Xs$ at $y_j$.  The
last two are invariant under $\mathrm{Gal}(\Qbar/\Q)$ because the maps are defined
over $\Q$; the first is $3$, $2$ or $1$ according as $j(y_j)$ is $0$, $1728$ or
neither, and these values are rational.  The function $G$ in the proof of Theorem~\ref{thm:elliptic} lies in
$\Q(X_0(N))$ and, by that proof, is constant on the fibre of $X_0(N)\to\Xs$ over
each $x_j$; so the numbers
$n_F(\tau_j)$ are the conjugates of the algebraic integer $n_F(\tau_1)$, and their
product is an integer of absolute value $\prod_j\nu_F(\tau_j)^{6\sz(N)}<1$, hence
zero.  Then one, and so each, $n_F(\tau_j)$ vanishes, that is
$F^{[e]}(\tau_j)=0$; by \eqref{eq:phi-ell} and linearity this is the claim.
\end{proof}

\begin{proposition}[Polynomials in the forms]\label{prop:poly}
Let $P\in\Z[X_1,\dots,X_g]$ be homogeneous of degree $m$, and let $x$, $h$,
$x_j$, $\tau_j$ be as in Corollary~\ref{cor:vanish}, with common order $e$.  Put
\[
\nu_P(\tau)=\Bigl(\frac{\sqrt N\,\im\tau}{T(\tau)}\Bigr)^{em}\,
\bigl\lvert P\bigl(f_1^{[e]}(\tau),\dots,f_g^{[e]}(\tau)\bigr)\bigr\rvert .
\]
Then $\nu_P(\tau_1)^{6\sz(N)}$ is the absolute value of an algebraic integer in
$\Q(x)$ whose conjugates have absolute values $\nu_P(\tau_j)^{6\sz(N)}$.  In
particular, if $\prod_{j}\nu_P(\tau_j)<1$, or if $\nu_P(\tau_1)M^{h-1}<1$ with
$M\ge1$ a bound for $\nu_P$ on a finite Galois-stable set of CM points
containing $x$, then $P$ vanishes at $\varphi(x_j)$ for every $j$.
\end{proposition}
\begin{proof}
Let $H_i=12^{e-1}\vartheta_{2e-2}\circ\cdots\circ\vartheta_2f_i$, cusp-integral of
weight $2e$ with $H_i(\tau_j)=f_i^{[e]}(\tau_j)$, as in the proof of
Theorem~\ref{thm:elliptic}.  Products of cusp-integral forms are cusp-integral
(the weights and the powers of $Q$ add), as are integral combinations of
cusp-integral forms of equal weight; so $H=P(H_1,\dots,H_g)$ is cusp-integral of
weight $2em$.  Theorem~\ref{thm:integral} with $k=2em$, $a=6\sz(N)$, $b=em$
shows that $G=H^a/D_N^{\,em}$ takes an algebraic integer value at $\tau_j$, of
absolute value $\nu_P(\tau_j)^{6\sz(N)}$.  For an Atkin--Lehner matrix $w$ the
proof of Theorem~\ref{thm:elliptic} gives
$f_i^{[e]}(w\tau_j)=\rho(\tau_j)^ef_i^{[e]}(\tau_j)$, hence
$H(w\tau_j)=\rho(\tau_j)^{em}H(\tau_j)$ by homogeneity, while
$D_N(w\tau_j)^{em}=\rho(\tau_j)^{6\sz(N)em}D_N(\tau_j)^{em}$.  So $G$ is constant
on the fibres of $X_0(N)\to\Xs$ over the $x_j$, and the argument of
Corollary~\ref{cor:vanish} applies.  Finally $P(\varphi(x_j))=0$ if and only if
$P(f_1^{[e]}(\tau_j),\dots,f_g^{[e]}(\tau_j))=0$, by \eqref{eq:phi-ell}.
\end{proof}

\begin{remark}[Use]\label{rem:use}
To prove that a rational point $P$ of the canonical model equals $\varphi(x)$
for a rational CM point $x$, apply Corollary~\ref{cor:vanish} with $h=1$ to
$g-1$ independent integral linear forms vanishing at $P$.  To prove that a
rational hyperplane contains a conjugate pair of quadratic CM points, apply it
with $h=2$; the pairing of the points need not be known, only a bound $M$ over
all CM points of the same discriminant.  To prove that a Galois-stable finite
set $Z$ of points of the canonical model, given by homogeneous equations, is the
set of images of the CM points of a given discriminant, apply
Proposition~\ref{prop:poly} to each equation of $Z$ and compare cardinalities.
Rationality of $x$, or its degree,
is an input from the theory of complex multiplication.  The quantity $T(\tau)$ is
computed from $\etanorm{d\tau}=\etanorm{z_d}$, with $z_d$ the root of the reduced
binary quadratic form equivalent to that of $d\tau$; there
$\lvert e^{2\pi iz_d}\rvert\le e^{-\pi\sqrt3}$, and
$\lvert\prod_n(1-q^n)\rvert\ge1-\lvert q\rvert/(1-\lvert q\rvert)$ gives a lower
bound for $T$ from one exponential.
\end{remark}

\subsection*{Genus two}
When $g=2$ the curve $\Xs$ is hyperelliptic and the canonical map is not an
embedding.  Write $S=\Q f\oplus\Q h$ with $f=q+\cdots$ and $h=q^2+\cdots$ up to
scalars, and put
\[
x=f/h,\qquad W=h\,\theta f-f\,\theta h,\qquad y=W/h^3=\theta(x)/h .
\]
Then $x$ and $y$ are functions on $\Xs$ and $y^2=P(x)$ for a polynomial
$P\in\Q[x]$ of degree $5$ or $6$; this is an identity between modular forms of
weight $12$, $W^2=\sum_kc_kf^kh^{6-k}$, and is proved by comparing Fourier
coefficients up to the Sturm bound.  The model lies in the weighted plane with
weights $(1,3,1)$, and the points with $h=0$ are $[1:W/f^3:0]$.

\begin{proposition}[Genus two]\label{prop:genus2}
Let $g=2$, let $\tau_0$ be a point whose stabilizer has order $e$, with image
$x_0\in\Xs$, and write $F_k=\theta^kF(\tau_0)$.  Put
\[
W^{\circ}=\bigl((e-1)!\bigr)^3\sum_{\substack{i+j=3e-2\\ e-1\le i,j\le 2e-1}}
\frac{j-i}{i!\,j!}\;h_i\,f_j .
\]
Then the point of the model corresponding to $x_0$ is
\[
[\,f_{e-1}:W^{\circ}:h_{e-1}\,]\ \text{ if } h_{e-1}\ne0,
\qquad
[\,1:W^{\circ}/f_{e-1}^{3}:0\,]\ \text{ if } h_{e-1}=0 .
\]
\end{proposition}
\begin{proof}
The first and third coordinates are given by \eqref{eq:phi-ell}; in particular
$f_k=h_k=0$ for $k<e-1$ and $(f_{e-1},h_{e-1})\ne(0,0)$.  Let
$v=2\pi i(\tau-\tau_0)$, so that $F=\sum_kF_kv^k/k!$ and $\theta=d/dv$.  The
coefficient of $v^n$ in $W=h\,f'-f\,h'$ is
\[
\sum_{i+j=n+1}(j-i)\,\frac{h_i}{i!}\,\frac{f_j}{j!},
\]
a sum over $i,j\ge e-1$.  Suppose $h_{e-1}\ne0$.  The function $y=W/h^3$ is
finite at $x_0$, because $x(x_0)$ is finite and $y^2=P(x)$.  So $W$ vanishes
at $\tau_0$ to order at least $3(e-1)$, the order of $h^3$, and $y(x_0)$ is
the ratio of the coefficients of $v^{3(e-1)}$ in $W$ and in $h^3$.  The latter
is $h_{e-1}^3/((e-1)!)^3$, and the former is the displayed sum with $n+1=3e-2$,
in which $i,j\le 2e-1$.  If $h_{e-1}=0$ the same argument applies to $W/f^3$,
which is the second coordinate in the chart $X=1$ and is finite on $Z=0$.
\end{proof}

\begin{remark}\label{rem:genus2}
To prove that a rational point $[X:Y:Z]$ of the model, with $X,Z$ coprime
integers, is the image of a rational CM point $x_0$: apply
Corollary~\ref{cor:vanish} with $h=1$ to $F=Zf-Xh$, and check by an enclosure
that $h_{e-1}\ne0$ (or $f_{e-1}\ne0$ when $Z=0$).  This gives $x(x_0)=X/Z$,
hence $y(x_0)=\pm Y/Z^3$.  When $Y\ne0$ the sign is settled by an enclosure
showing $Z^3W^{\circ}+Y\,h_{e-1}^3\ne0$ (respectively
$X^3W^{\circ}+Y\,f_{e-1}^3\ne0$).  No integrality statement for $W$ is needed.
\end{remark}

\section{An explicit bound for Fourier coefficients}\label{sec:tail}

A rigorous upper bound for $\lvert F(\tau)\rvert$ from a truncated $q$-expansion
needs a bound on the coefficients of $F$.

\begin{lemma}\label{lem:star-basis}
The space $S$ has a basis consisting of the forms
\[
h_\phi(\tau)=\sum_{d\mid N/M}d\,\phi(d\tau),
\]
where $M$ runs over the divisors of $N$ and $\phi$ over the normalized newforms
of level $M$ all of whose Atkin--Lehner eigenvalues are $+1$.  These are
simultaneous eigenforms of the Hecke operators $T_\ell$, $\ell\nmid N$, with
pairwise distinct systems of eigenvalues, and
\[
\lvert a_n(h_\phi)\rvert\le\kappa_N\,\sz(n)\sqrt n,\qquad
\kappa_N=\prod_{p\mid N}(1+\sqrt p).
\]
\end{lemma}
\begin{proof}
By \cite[Thm.~5]{AL}, $S_2(\Gamma_0(N))$ is the direct sum over $M\mid N$ and
newforms $\phi$ of level $M$ of the spaces spanned by
$u_d=\phi|_2\smat d001=d\,\phi(d\tau)$, $d\mid N/M$.  Let $p\mid N$ and let
$w_p=\smat{px}{y}{Nz}{pt}$.  If $p\mid M$, then
$\smat d001w_p\smat d001^{-1}=\smat{px}{dy}{Nz/d}{pt}$ is an Atkin--Lehner matrix
of level $M$, so $u_d|w_p=\varepsilon_p(\phi)u_d$.  If $p\mid N/M$ and
$p\nmid d$, then $\smat d001w_p\smat{dp}001^{-1}=\smat{x}{dy}{Nz/(dp)}{pt}$ lies
in $\Gamma_0(M)$, so $u_d|w_p=u_{dp}$, and $u_{dp}|w_p=u_d$ because $w_p^2$ acts trivially.  Hence the
$W(N)$-invariants in the span of the $u_d$ are zero unless all
$\varepsilon_p(\phi)=1$, in which case they are spanned by $\sum_du_d=h_\phi$.
Distinctness of the eigenvalue systems is strong multiplicity one \cite[Thm.~4]{AL}.

For the bound, $\lvert a_\ell(\phi)\rvert\le2\sqrt\ell$ for $\ell\nmid M$ by
Eichler--Shimura and Weil \cite[Ch.~7]{Shimura} (a case of \cite{Deligne}), and
$a_p(\phi)=\pm1$ for $p\mid M$ \cite[Thm.~3]{AL}; by multiplicativity
$\lvert a_n(\phi)\rvert\le\sz(n)\sqrt n$.  Then
$\lvert a_n(h_\phi)\rvert\le\sum_{d\mid(n,N/M)}d\,\sz(n/d)\sqrt{n/d}
\le\sz(n)\sqrt n\sum_{d\mid N/M}\sqrt d$.
\end{proof}

\begin{proposition}\label{prop:tail}
Let $\mathcal T$ be an integral linear combination of Hecke operators $T_\ell$,
$\ell\nmid N$, whose characteristic polynomial on $S$ is squarefree.  For
$F\in S$ with rational coefficients put $u_j=a_1(\mathcal T^jF)$ and
$\mathcal M_{jl}=\operatorname{tr}(\mathcal T^{j+l}|S)$ for $0\le j,l\le g-1$.
Then $\mathcal M$ is positive definite and
\[
\lvert a_n(F)\rvert\le C_F\,\sz(n)\sqrt n,\qquad
C_F=\kappa_N\sqrt{g\cdot u^{\mathsf T}\mathcal M^{-1}u}\;.
\]
Consequently, for $\lvert q\rvert=r<1$,
\[
\Bigl\lvert F(\tau)-\sum_{n\le m}a_n(F)q^n\Bigr\rvert
\le2C_F\,\frac{r^{m+1}\,(m+1-mr)}{(1-r)^2}.
\]
\end{proposition}
\begin{proof}
Write $F=\sum_\phi c_\phi h_\phi$ over the basis of Lemma~\ref{lem:star-basis};
the eigenvalues $\lambda_\phi$ of $\mathcal T$ are real and distinct, and the
$c_\phi$ are real because $F$ is rational and the $h_\phi$ are permuted by
$\mathrm{Gal}(\Qbar/\Q)$ with totally real coefficients.  Since $a_1(h_\phi)=1$,
$u_j=\sum_\phi c_\phi\lambda_\phi^j$, that is $u=V^{\mathsf T}c$ with
$V=(\lambda_\phi^j)$ an invertible Vandermonde matrix, and
$\mathcal M=V^{\mathsf T}V$.  Hence
$u^{\mathsf T}\mathcal M^{-1}u=\sum_\phi c_\phi^2$, and
$\sum_\phi\lvert c_\phi\rvert\le\sqrt{g\sum c_\phi^2}$.  The tail bound uses
$\sz(n)\le2\sqrt n$ and $\sum_{n>m}nr^n=r^{m+1}(m+1-mr)/(1-r)^2$.
\end{proof}

All quantities in Proposition~\ref{prop:tail} are computed in exact rational
arithmetic from the $q$-expansions of a basis of $S$.

\section{Worked examples at level \texorpdfstring{$67$}{67}}\label{sec:ex}

The smallest squarefree level for which $\Xs$ has genus at least $2$ is $N=67$,
where $g=2$ and $\sz(N)=2$.  Three statements are proved here with every number
displayed: a rational CM point, an elliptic point, and a conjugate pair of
quadratic CM points.  The only inputs are the first $80$ Fourier coefficients of
two forms; the arithmetic is that of a pocket calculator, although nobody would
want to sum the series without a machine.  Decimals are rounded in the safe
direction.

\subsection*{Setup}
The space $S$ has the basis
\begin{align*}
f&=q-3q^3-3q^4-3q^5+q^6+4q^7+3q^8+5q^9-3q^{11}+\cdots,\\
h&=q^2-q^3-3q^4+3q^7+4q^8+3q^9-3q^{10}-2q^{11}+\cdots,
\end{align*}
and with $x=f/h$, $y=(h\,\theta f-f\,\theta h)/h^3$ the model is
\[
y^2=P(x)=x^6-4x^5+6x^4-6x^3+9x^2-14x+9 .
\]

\emph{Coefficient bounds.}  From $a_n(T_2F)=a_{2n}(F)+2a_{n/2}(F)$ and the
coefficients above, $T_2f=-h$ and $T_2h=f-3h$.  The characteristic polynomial
$X^2+3X+1$ is squarefree, so Proposition~\ref{prop:tail} applies with
$\mathcal T=T_2$:
\[
\mathcal M=\begin{pmatrix}2&-3\\-3&7\end{pmatrix},\qquad
\mathcal M^{-1}=\frac15\begin{pmatrix}7&3\\3&2\end{pmatrix},\qquad
\kappa_{67}=1+\sqrt{67}\le9.186 .
\]
For $F=\alpha f+\beta h$ we have $u=(a_1(F),a_2(F))=(\alpha,\beta)$ and
$C_F=\kappa_{67}\sqrt{2(7\alpha^2+6\alpha\beta+2\beta^2)/5}$, so
\[
C_f\le15.38,\qquad C_h\le8.22,\qquad C_{f-2h}\le10.07,\qquad C_{f+h}\le22.5 .
\]
For derivatives, $n^k\lvert a_n(F)\rvert\le2C_Fn^{k+1}$, and the terms
$n^{k+1}r^n$, $n>m$, decrease at least geometrically with ratio
$\rho_k=r\bigl(\tfrac{m+2}{m+1}\bigr)^{k+1}$, so for $\lvert q\rvert=r$ and $\rho_k<1$
\begin{equation}\label{eq:tail-theta}
\Bigl\lvert\theta^kF(\tau)-\sum_{n\le m}n^ka_n(F)q^n\Bigr\rvert
\le\frac{2C_F\,(m+1)^{k+1}r^{m+1}}{1-\rho_k}.
\end{equation}

\emph{The factor $T$.}  Let $D<0$ be a discriminant of class number one and
$\tau=(-B+\sqrt D)/134$ a root of a form $[67,B,C]$ of discriminant $D$.  The
forms of $\tau$ and of $67\tau$, namely $[67,B,C]$ and $[1,B,67C]$, are both
equivalent to the principal form, so
$\etanorm{\tau}=\etanorm{67\tau}=\etanorm{z_D}$ with
$\im z_D=\sqrt{\lvert D\rvert}/2$, and
\[
T(\tau)=\etanorm{z_D}^4=\tfrac12\sqrt{\lvert D\rvert}\;e^{-\pi\sqrt{\lvert D\rvert}/6}\,
\Bigl\lvert\prod_{n\ge1}(1-q_z^n)\Bigr\rvert^4,
\qquad\lvert q_z\rvert=e^{-\pi\sqrt{\lvert D\rvert}} .
\]  Since
$\sqrt{67}\,\im\tau=\sqrt{\lvert D\rvert}/(2\sqrt{67})$, Remark~\ref{rem:use}
gives, for $\lvert D\rvert\ge43$,
\begin{equation}\label{eq:scale67}
\frac{\sqrt{67}\,\im\tau}{T(\tau)}\le(1+10^{-8})\,\frac{e^{\pi\sqrt{\lvert D\rvert}/6}}{\sqrt{67}} .
\end{equation}

\begin{example}[A rational CM point, $D=-43$]\label{ex:43}
The class number of $D=-43$ is one and $15^2+43=4\cdot67$, so
$\tau=(15+\sqrt{-43})/134$ is a Heegner point with rational image
$x_0\in X_0(67)^*(\Q)$.  Its stabilizer is trivial: an integer matrix fixing $\tau$ is
$\smat{(t+15u)/2}{-u}{67u}{(t-15u)/2}$, of determinant $(t^2+43u^2)/4$, and
the only solutions of $t^2+43u^2=4Q$ with $Q\in\{1,67\}$, $u\ge1$ are
$(t,u)=(\pm15,1)$ with $Q=67$, for which the diagonal entries are $0$ and
$\pm15$, not both divisible by $67$.  We claim that $x_0$ is the point $(2:-1:1)$
of the model.

Here $q=r\,e^{2\pi i\cdot15/134}$ with $r=e^{-\pi\sqrt{43}/67}\le0.7354$, and by
\eqref{eq:scale67} the factor in $\nu$ is at most $3.7853$.  Take $F=f-2h$.
With $m=35$ terms the partial sum has absolute value at most $2.8\cdot10^{-4}$,
and the tail bound of Proposition~\ref{prop:tail} is
\[
2\cdot10.07\cdot\frac{r^{36}(36-35r)}{(1-r)^2}\le0.047 .
\]
Hence $\nu_F(\tau)\le3.7853\cdot0.0473<0.18<1$, and Theorem~\ref{thm:main-intro}
gives $f(\tau)=2h(\tau)$ exactly.  With $60$ terms,
\[
h(\tau)=1.58730-0.69391\,i,\qquad
\theta f(\tau)=13.8795-8.3559\,i,\qquad
\theta h(\tau)=7.9587-5.2794\,i,
\]
each with error at most $4\cdot10^{-3}$ by \eqref{eq:tail-theta}.  So
$h(\tau)\ne0$ and $x(x_0)=2$.  Since $P(2)=1$ we get $y(x_0)=\pm1$, that is
$W(\tau)=\pm h(\tau)^3$ for $W=h\,\theta f-f\,\theta h$.  Numerically
$W(\tau)=-1.706+4.911\,i$ and $h(\tau)^3=1.706-4.911\,i$, with errors at most
$0.05$, so $\lvert W(\tau)-h(\tau)^3\rvert\ge10$ and the sign is $-1$.

The integers of Theorem~\ref{thm:main-intro} can be seen as well.  Summing more
terms gives $\nu_h(\tau)=6.557438\ldots$, which agrees with $\sqrt{43}$ to all
the digits computed, and $\nu_f=2\nu_h$; so
\[
n_h(\tau)=43^6,\qquad n_f(\tau)=2^{12}\cdot43^6 .
\]
Corollary~\ref{cor:height-intro} alone also proves $x(x_0)=2$: it gives
$\varphi(x_0)=(a_1:a_2)$ with $\lvert a_1\rvert\le13$ and $\lvert a_2\rvert\le6$,
where $a_2\ne0$ because $h(\tau)\ne0$; a fraction $a_1/a_2\ne2$ with
$\lvert a_2\rvert\le6$ differs from $2$ by at least $1/6$, whereas
$\lvert f(\tau)/h(\tau)-2\rvert\le0.0473/1.72<1/6$.
\end{example}

\begin{example}[An elliptic point, $D=-67$]\label{ex:67}
Let $\tau_0=\tfrac12+\tfrac{i}{2\sqrt{67}}$, the root of $[67,-67,17]$, of
discriminant $-67$ and class number one.  It is fixed by the Atkin--Lehner
matrix $\smat{67}{-34}{134}{-67}$ of determinant $67$, and the stabilizer has
order $e=2$ (the only solution of $t^2+67u^2=4Q$ with $Q\in\{1,67\}$, $u\ge1$
is $(t,u)=(0,2)$, which gives this matrix).  So $f(\tau_0)=h(\tau_0)=0$, and the image $x_0$ is described by
first derivatives.  We claim that $x_0=(-1:7:1)$.

Now $q=-e^{-\pi/\sqrt{67}}$ is real, $r\le0.68127$, and all values below are
real.  By \eqref{eq:scale67} and Theorem~\ref{thm:elliptic},
\[
\Bigl(\frac{\sqrt{67}\,\im\tau_0}{T(\tau_0)}\Bigr)^2\le78.807,
\qquad
\nu_F(\tau_0)\le78.807\cdot12\cdot\lvert\theta F(\tau_0)\rvert .
\]
Take $F=f+h$.  With $m=60$ terms the partial sum of $\theta F$ is at most
$1.5\cdot10^{-7}$ in absolute value and the bound \eqref{eq:tail-theta} with
$k=1$ is at most $3.9\cdot10^{-5}$, so $\nu_F(\tau_0)<0.04<1$ and
$\theta f(\tau_0)=-\theta h(\tau_0)$.  With $80$ terms, writing
$F_k=\theta^kF(\tau_0)$,
\[
f_1=1.7003788,\quad h_1=-1.7003788,\quad
f_3=149.0693,\quad h_3=-88.3522,
\]
with errors at most $2\cdot10^{-4}$.  So $h_1\ne0$ and $x(x_0)=f_1/h_1=-1$ by
\eqref{eq:phi-ell}.  For $e=2$ the sum in Proposition~\ref{prop:genus2} has the
terms $(i,j)=(1,3),(2,2),(3,1)$, and
\[
W^\circ=\tfrac13\,(h_1f_3-h_3f_1)=-34.414,\qquad h_1^3=-4.9163,
\qquad W^\circ/h_1^3=7.000 .
\]
Since $P(-1)=49$, the value $y(x_0)=W^\circ/h_1^3$ is $\pm7$, and
$W^\circ+7h_1^3=-68.8\ne0$ rules out $-7$.

Here $\nu_f(\tau_0)=78.806\cdot12\cdot1.7003788=1608.00$, which is $24\cdot67$
to the digits computed; so the integer of Theorem~\ref{thm:elliptic} is
$n_f(\tau_0)=(24\cdot67)^{12}$.
\end{example}

\begin{example}[A conjugate pair, $D=-115$]\label{ex:115}
The class number of $D=-115$ is two, and $35^2+115=4\cdot67\cdot5$.  The CM
points of discriminant $-115$ on $X_0(67)^*$ are the images $x_\pm$ of
$\tau_\pm=(\pm35+\sqrt{-115})/134$, the roots of $[67,\mp35,5]$; they form a
Galois-stable set.  We claim that $x(x_\pm)$ are the roots of $2x^2-3x+4$,
that is $(3\pm\sqrt{-23})/4$.

The form of $67\tau_\pm$ is $[1,\mp35,335]$, equivalent to the principal form
$[1,1,29]$, and $[67,\mp35,5]$ is equivalent to the reduced form $[5,\pm5,7]$.
With one exponential each, as in Remark~\ref{rem:use},
\begin{align*}
\etanorm{67\tau_\pm}^4&\ge\tfrac{\sqrt{115}}2\,e^{-\pi\sqrt{115}/6}(1-10^{-9})\ge0.019532,\\
\etanorm{\tau_\pm}^4&\ge\tfrac{\sqrt{115}}{10}\,e^{-\pi\sqrt{115}/30}\Bigl(1-\tfrac{\rho}{1-\rho}\Bigr)^4\ge0.34719,
\end{align*}
where $\rho=e^{-\pi\sqrt{115}/5}\le0.0011851$.  Hence
$T(\tau_\pm)\ge0.08234$, and
\[
\sqrt{67}\,\im\tau_\pm=\frac{\sqrt{115}}{2\sqrt{67}}\le0.65506,
\qquad
\Bigl(\frac{\sqrt{67}\,\im\tau_\pm}{T(\tau_\pm)}\Bigr)^2\le63.3 .
\]
Apply Proposition~\ref{prop:poly} to $G=2X_1^2-3X_1X_2+4X_2^2$, of degree
$m=2$, with $e=1$: $\nu_G(\tau)\le63.3\,\lvert G(f(\tau),h(\tau))\rvert$.
Here $r=e^{-\pi\sqrt{115}/67}\le0.60482$, and with $40$ terms the partial sums
are
\[
S_f=-0.440708\pm0.873544\,i,\qquad S_h=-0.688937\pm0.063384\,i,
\]
with tails at most $\varepsilon=3.7\cdot10^{-6}$.  Then
$\lvert G(S_f,S_h)\rvert\le2.5\cdot10^{-8}$.  Since
$\lvert S_f\rvert$ and $\lvert S_h\rvert$ are at most $1$, each of the three monomials changes by
at most $2\varepsilon+\varepsilon^2$ when the partial sums are replaced by the
true values; the coefficients of $G$ have absolute sum $9$, so
\[
\lvert G(f(\tau_\pm),h(\tau_\pm))\rvert\le2.5\cdot10^{-8}+9(2\varepsilon+\varepsilon^2)\le6.7\cdot10^{-5},
\qquad\nu_G(\tau_\pm)\le0.0043 .
\]
If $x_+$ and $x_-$ are conjugate then $\nu_G(\tau_+)\nu_G(\tau_-)<1$, and if
each is rational then each $\nu_G<1$; in both cases $G$ vanishes at
$\varphi(x_\pm)=(f(\tau_\pm):h(\tau_\pm))$.  As $h(\tau_\pm)\ne0$, the values
$x(x_\pm)=S_f/S_h+O(10^{-5})=0.75\mp1.19896\,i$ are roots of $2x^2-3x+4$, one
each, so the points are conjugate over $\Q(\sqrt{-23})$.
\end{example}

The routines in \texttt{src/cm\_certify.m} carry out these steps with rational
balls in place of rounded decimals.

\section{Numerical checks}\label{sec:num}

The statements were tested in Magma on the $3000$-term expansions of integral
bases of $S$; these are consistency checks, not part of the proofs.
The distance to the nearest integer is limited by the truncation of the series.

\begin{center}
\begin{tabular}{llp{7.6cm}}
$N$ & points tested & outcome\\\hline
$137$ & $7$ rational CM points (one with $e=2$) & $n_{f_i}$ within $10^{-43}$ of an integer, or better\\
$154$ & $10$ rational CM points (two with $e=2$) & within $10^{-13}$, or better\\
$330$ & $6$ rational CM points & $4$ integral to $10^{-13}$ or better; $D=-120,-24$ inconclusive (integers of $147$--$181$ digits exceed the truncation accuracy)\\
$311$ & conjugate pairs, $D=-123,-232$ & trace and norm of $n_{f_i}$ integral to $10^{-124}$\\
\end{tabular}
\end{center}

The table used the equivalent normalization
$(F^{[e]})^{6\sz(N)/e}/D_N$ at elliptic points; only $e\le2$ occurred.  At
$N=137,154,311,330,910$ the bounds of Corollary~\ref{cor:height} at non-elliptic
rational CM points ranged from $2.6$ to $296$, the constants $C_{f_i}$ of
Proposition~\ref{prop:tail} from $12$ to $306$, and between $20$ and $410$ terms
of the $q$-expansion sufficed for the tail bound, in the normalization of $\nu$, to fall below $1/4$.

\begin{thebibliography}{99}
\bibitem{AL} A.~O.~L. Atkin and J.~Lehner, \emph{Hecke operators on $\Gamma_0(m)$},
Math. Ann. \textbf{185} (1970), 134--160.
\bibitem{BruinierOno} J.~H. Bruinier and K.~Ono, \emph{Algebraic formulas for the
coefficients of half-integral weight harmonic weak Maass forms}, Adv. Math.
\textbf{246} (2013), 198--219; arXiv:1104.1182.
\bibitem{Cox} D.~A. Cox, \emph{Primes of the form $x^2+ny^2$}, 2nd ed., Wiley, 2013.
\bibitem{Deligne} P.~Deligne, \emph{La conjecture de Weil.~I}, Publ. Math. IH\'ES
\textbf{43} (1974), 273--307.
\bibitem{DS} F.~Diamond and J.~Shurman, \emph{A first course in modular forms},
Graduate Texts in Math. 228, Springer, 2005.
\bibitem{GrossZagier} B.~H. Gross and D.~B. Zagier, \emph{On singular moduli},
J. reine angew. Math. \textbf{355} (1985), 191--220.
\bibitem{KubertLang} D.~S. Kubert and S.~Lang, \emph{Modular units}, Grundlehren
der math. Wissenschaften 244, Springer, 1981.
\bibitem{Lang} S.~Lang, \emph{Elliptic functions}, 2nd ed., Graduate Texts in
Math. 112, Springer, 1987.
\bibitem{Shimura} G.~Shimura, \emph{Introduction to the arithmetic theory of
automorphic functions}, Princeton Univ. Press, 1971.
\bibitem{SilvermanAT} J.~H. Silverman, \emph{Advanced topics in the arithmetic of
elliptic curves}, Graduate Texts in Math. 151, Springer, 1994.
\bibitem{Zagier} D.~Zagier, \emph{Elliptic modular forms and their applications},
in: The 1-2-3 of modular forms, Universitext, Springer, 2008, 1--103.
\end{thebibliography}
\end{document}
